\begin{recomendations}
El desarrollo de la Base de Imágenes de DermaUH ha sentado las bases fundacionales de una infraestructura 
de datos clínicos que, por su naturaleza, admite y requiere un proceso continuo de evolución y 
enriquecimiento. A partir de la experiencia acumulada durante la investigación y el análisis de 
las tendencias actuales en dermatología computacional, se formulan las siguientes recomendaciones 
para las etapas sucesivas del proyecto:

\begin{enumerate}
    \item \textbf{Incorporación de máscaras de segmentación.} Se recomienda extender el esquema de datos 
    para almacenar máscaras de segmentación binarias asociadas a cada imagen dermatoscópica, siguiendo la 
    metodología establecida por iniciativas como IMA++. La inclusión de anotaciones de segmentación a 
    nivel de píxel habilitaría el entrenamiento de arquitecturas de redes neuronales para la delimitación 
    automática de lesiones, ampliando significativamente el espectro de aplicaciones de inteligencia 
    artificial soportadas por el repositorio.

    \item \textbf{Implementación de técnicas de aumento de datos y generación sintética.} Dada la 
    previsible escasez inicial de imágenes en determinadas categorías diagnósticas y fototipos de piel, 
    se recomienda explorar la integración de técnicas de aumento de datos y modelos generativos para la 
    síntesis de imágenes dermatológicas, siguiendo la línea de investigación establecida por herramientas 
    como DermaSynth, con el objetivo de mitigar el desbalance estadístico inherente a las fases tempranas 
    de la recolección.

    \item \textbf{Desarrollo de un módulo de control de calidad de imagen.} Se sugiere incorporar filtros 
    automáticos de calidad de imagen basados en redes neuronales ligeras que evalúen indicadores de 
    nitidez, exposición y enfoque en el momento de la carga, inspirados en el \textit{framework} DermAI. 
    Este componente elevaría la homogeneidad técnica del \textit{corpus} y reduciría la necesidad de 
    curación manual posterior.

    \item \textbf{Expansión del cuadro de mando analítico.} Se recomienda enriquecer el módulo de 
    estadísticas con visualizaciones temporales que permitan analizar tendencias epidemiológicas a lo 
    largo del tiempo, así como incorporar tabulaciones cruzadas más complejas que combinen variables 
    demográficas, clínicas y geográficas para habilitar estudios de epidemiología descriptiva de mayor 
    profundidad.

    \item \textbf{Contribución al ecosistema ISIC.} Una vez alcanzado un volumen significativo de 
    imágenes con metadatos validados y consentimiento informado verificado, se recomienda explorar la 
    contribución de los datos cubanos al archivo ISIC, aprovechando la compatibilidad estructural del 
    esquema de metadatos diseñado. Esta acción posicionaría a la Universidad de La Habana como un centro 
    contribuyente reconocido en la comunidad internacional de dermatología computacional y ampliaría la 
    representatividad fenotípica de los repositorios globales.

    \item \textbf{Publicación científica de los resultados.} Se recomienda elaborar publicaciones 
    científicas que documenten tanto la arquitectura del sistema como los primeros análisis 
    epidemiológicos derivados de los datos recopilados, dirigidas a revistas especializadas en 
    informática médica y dermatología digital. La originalidad del enfoque ---un repositorio soberano 
    con datos de una población fenotípicamente diversa y subrepresentada--- constituye un aporte de alto 
    interés para la comunidad científica internacional.
\end{enumerate}

La ejecución de estas recomendaciones permitirá evolucionar la Base de Imágenes de DermaUH desde un 
repositorio institucional operativo hacia una plataforma de referencia para la investigación dermatológica 
computacional en el contexto latinoamericano y caribeño, consolidando el liderazgo de la Universidad de 
La Habana en el desarrollo de tecnologías de salud soberanas y equitativas.
\end{recomendations}
